\documentclass[]{article}
\usepackage{lmodern}
\usepackage{amssymb,amsmath}
\usepackage{ifxetex,ifluatex}
\usepackage{fixltx2e} % provides \textsubscript
\ifnum 0\ifxetex 1\fi\ifluatex 1\fi=0 % if pdftex
  \usepackage[T1]{fontenc}
  \usepackage[utf8]{inputenc}
\else % if luatex or xelatex
  \ifxetex
    \usepackage{mathspec}
  \else
    \usepackage{fontspec}
  \fi
  \defaultfontfeatures{Ligatures=TeX,Scale=MatchLowercase}
\fi
% use upquote if available, for straight quotes in verbatim environments
\IfFileExists{upquote.sty}{\usepackage{upquote}}{}
% use microtype if available
\IfFileExists{microtype.sty}{%
\usepackage{microtype}
\UseMicrotypeSet[protrusion]{basicmath} % disable protrusion for tt fonts
}{}
\usepackage[margin=1in]{geometry}
\usepackage{hyperref}
\PassOptionsToPackage{usenames,dvipsnames}{color} % color is loaded by hyperref
\hypersetup{unicode=true,
            pdftitle={Master thesis},
            pdfauthor={Sylvain SCHMITT},
            colorlinks=true,
            linkcolor=Maroon,
            citecolor=Blue,
            urlcolor=Blue,
            breaklinks=true}
\urlstyle{same}  % don't use monospace font for urls
\usepackage{natbib}
\bibliographystyle{plainnat}
\usepackage{longtable,booktabs}
\usepackage{graphicx,grffile}
\makeatletter
\def\maxwidth{\ifdim\Gin@nat@width>\linewidth\linewidth\else\Gin@nat@width\fi}
\def\maxheight{\ifdim\Gin@nat@height>\textheight\textheight\else\Gin@nat@height\fi}
\makeatother
% Scale images if necessary, so that they will not overflow the page
% margins by default, and it is still possible to overwrite the defaults
% using explicit options in \includegraphics[width, height, ...]{}
\setkeys{Gin}{width=\maxwidth,height=\maxheight,keepaspectratio}
\IfFileExists{parskip.sty}{%
\usepackage{parskip}
}{% else
\setlength{\parindent}{0pt}
\setlength{\parskip}{6pt plus 2pt minus 1pt}
}
\setlength{\emergencystretch}{3em}  % prevent overfull lines
\providecommand{\tightlist}{%
  \setlength{\itemsep}{0pt}\setlength{\parskip}{0pt}}
\setcounter{secnumdepth}{5}
% Redefines (sub)paragraphs to behave more like sections
\ifx\paragraph\undefined\else
\let\oldparagraph\paragraph
\renewcommand{\paragraph}[1]{\oldparagraph{#1}\mbox{}}
\fi
\ifx\subparagraph\undefined\else
\let\oldsubparagraph\subparagraph
\renewcommand{\subparagraph}[1]{\oldsubparagraph{#1}\mbox{}}
\fi

%%% Use protect on footnotes to avoid problems with footnotes in titles
\let\rmarkdownfootnote\footnote%
\def\footnote{\protect\rmarkdownfootnote}

%%% Change title format to be more compact
\usepackage{titling}

% Create subtitle command for use in maketitle
\newcommand{\subtitle}[1]{
  \posttitle{
    \begin{center}\large#1\end{center}
    }
}

\setlength{\droptitle}{-2em}
  \title{Master thesis}
  \pretitle{\vspace{\droptitle}\centering\huge}
  \posttitle{\par}
  \author{Sylvain SCHMITT}
  \preauthor{\centering\large\emph}
  \postauthor{\par}
  \predate{\centering\large\emph}
  \postdate{\par}
  \date{2017-05-30}

\usepackage{booktabs}
\usepackage{amsthm}
\makeatletter
\def\thm@space@setup{%
  \thm@preskip=8pt plus 2pt minus 4pt
  \thm@postskip=\thm@preskip
}
\makeatother

\begin{document}
\maketitle

\date{}

\thispagestyle{empty}
\begin{center}
\includegraphics{images/main.pdf}
\end{center}

\setlength{\abovedisplayskip}{-5pt}
\setlength{\abovedisplayshortskip}{-5pt}

{
\hypersetup{linkcolor=black}
\setcounter{tocdepth}{2}
\tableofcontents
}
\section*{Preface}\label{preface}
\addcontentsline{toc}{section}{Preface}

This is the online html version of my master thesis. The pdf version is
available for download with the download button on top of the windows.

Have a nice reading.

\emph{Sylvain}

\section*{Résumé et Abstract}\label{resume-et-abstract}
\addcontentsline{toc}{section}{Résumé et Abstract}

Écrire le résumé français ici\ldots{}

Write the english abstract here\ldots{}

\section*{Acknowledgments}\label{acknowledgments}
\addcontentsline{toc}{section}{Acknowledgments}

I would like to thank\ldots{}

\section*{Introduction}\label{introduction}
\addcontentsline{toc}{section}{Introduction}

\begin{itemize}
\tightlist
\item
  Introduce why tropical rainforest are such a good study case
\item
  Define biodiversity
\item
  Define ecosystem functionning
\item
  Define ecosystem services
\item
  Introduce tropical sylviculture
\item
  Define selective logging
\item
  Introduce biodiversity erosion in tropical rainforests ant its impact
\item
  Introduce Loreau partitionning
\item
  Introduce the interest of modelling
\item
  Introduce the idea to look at biodiversity as both a parameter and an
  output
\end{itemize}

Tropical rainforests have fascinated ecologists due to their outstanding
diversity \citep{connell_diversity_1978}.

Tropical forests are primary ecosystems in terms of biodiversit and
carbon storage \citep{Lewis2004}.

High biodiversity from tropical rainforests is the source of many
ecosystem functions, which support ecosystem services.

Biodiversity and ecosystem functionning field is emmering
\citep{Loreau2000}.

Selective logging is increasing in tropical forests.

Tropical logging accounts for one eight of global timber production
\citep{Blaser2011}.

Selective logging represents a major threat to biodiversity
\citep{Gibson2013}.

\section{Materials and methods}\label{materials-and-methods}

\subsection{TROLL simulator}\label{troll-simulator}

\subsubsection{Abiotic environment}\label{abiotic-environment}

\subsubsection{Photosynthesis}\label{photosynthesis}

\subsubsection{Autotrpohic respiration}\label{autotrpohic-respiration}

\subsubsection{Carbon uptake}\label{carbon-uptake}

\subsubsection{Tree growth}\label{tree-growth}

\subsubsection{Seed dispersion, production and
recruitment}\label{seed-dispersion-production-and-recruitment}

\subsubsection{Mortality}\label{mortality}

\subsection{TROLL sensitivity
analysis}\label{troll-sensitivity-analysis}

\subsubsection{Functional traits}\label{functional-traits}

\subsubsection{Seed rain}\label{seed-rain}

\subsection{Disturbance}\label{disturbance}

\subsubsection{Disturbance module}\label{disturbance-module}

\subsubsection{Design of experiment}\label{design-of-experiment}

\subsubsection{Outputs anlaysis ?}\label{outputs-anlaysis}

\paragraph{Resistance and resilience
metrics}\label{resistance-and-resilience-metrics}

\paragraph{Biodiversity partitioning}\label{biodiversity-partitioning}

\subsection{Selective logging}\label{selective-logging}

\subsubsection{Selective logging module}\label{selective-logging-module}

\paragraph{Designation}\label{designation}

\paragraph{Selection}\label{selection}

\paragraph{Rotten trees}\label{rotten-trees}

\paragraph{Felling}\label{felling}

\paragraph{Tracks}\label{tracks}

\paragraph{Gap damages}\label{gap-damages}

\subsubsection{Design of experiment}\label{design-of-experiment-1}

\subsubsection{Outputs analysis ?}\label{outputs-analysis}

\paragraph{Resistance and resilience
metrics}\label{resistance-and-resilience-metrics-1}

\paragraph{Biodiversity partitioning}\label{biodiversity-partitioning-1}

\section{Results}\label{results}

\section{Discussion}\label{discussion}

\bibliography{bibliography/library.bib}


\end{document}
